% 致谢
% 原书页码：377；标题
\BookUnnumbered{致谢}
\par
% 原书页码：377；段落
书写到“致谢”大概都是五味杂陈的。写书的过程是人生一段难%
忘的旅程，往事种种涌上心头。本书是我的第一部中文专著，承载着%
我许多年轻时的梦想。从2004年出国留学，到2019年拿到博士学位，%
历时十五年之久。其间因为导师职务变动，因为英文不好，因为选题%
不被接受，因为考试不过关，因为不确定是不是能写出有意义的东西，%
我很多次游走在放弃的边缘。现在回想起来，我依然不知道为什么还%
能坚持一路走来。如果确有原因的话，我想可能是读大学时我真的有%
过做学问的梦想，内心深处始终有继续写作的召唤。学术研究难以用%
功利来衡量，我也深知不是每一份坚持都会有结果，正因如此，在这%
部书完稿时，我感慨良多。%
\par
% 原书页码：377–378；段落（跨页续段）
本书是在我博士论文以及先期出版的英文专著的基础上，扩充%
发展而来。2019年5月，我博士论文答辩后，委员会五位教授都鼓%
励我去联系出版社出版。我论文中有些不同于前人的新论断。答辩%
委员会中唯一的华人教授李弘祺老师私下跟我讲，如果我的论断是%
正确的，那么这一时段的很多历史都需要改写，最起码很多网上的%
历史词条需要重写。李老师专攻宋史，嘱我找出版社前，先联系跟%
我领域直接相关的专家，进一步核实我的论点。一位美国的教授和%
一位澳大利亚的教授分别给了我回信，二人正好是跟我研究最有关%
联的两位学者。他们都看出了我论文中原创的部分，分别从不同角%
度提出了他们的疑问，我也做了解答。两位学者表示我发现的大方%
向可能是正确的，鼓励我继续深挖，争取做出更精确的描述。特别%
是美国的教授，发给我了几份他早年在梵蒂冈查阅到的原始资料，%
这些材料正好是支持我观点的直接证据。他还提到一本1980年代推%
出的法文专著，专门把相关章节用手机拍了图片发给我，告诉我里%
面的研究就是我证据薄弱部分的支持。跟两位学者交流后，我多了%
几分底气，便开始着手联系出版社。%
\par
% 原书页码：378；段落
我在网上找到了几所美国大学出版社的电邮地址，开始发摘要%
投稿。第一轮投了六所大学的出版社，三所直接拒绝了我，三所要%
去了我的书稿全文进行审核。不过后来也都拒绝了我，说书稿博士%
论文味道太浓，声称他们必须要考虑出版图书的市场。如何去除书%
稿论文感其实是世界所有文科博士计划出书时都会面临的问题。不%
仅网上已经有过很多讨论和经验总结，芝加哥大学出版社还出有专%
书讲解，书名就叫 From Dissertation to Book（《从博士论文到书》）。我%
花了三个月来修改文字和结构，学着把“论文”变成“书”。这时期，%
我的导师 Richard Lufrano提醒我，跟天主教相关的内容，也许欧洲%
的出版社更有兴趣。%
\par
% 原书页码：379；段落
于是我想到了荷兰的博睿（Brill）出版社，我在写论文时就发现%
好多相关专著都是这里出版的。我把重新编好的书稿投到博睿后，学%
术编辑很快给我回复，告诉我已将书稿转给出版社责任编辑了。过了%
几天，责任编辑告诉我愿意受理我的书稿，外送匿名评审。在欧美，%
对于新毕业的博士来说，出版中很难的一步是出版社同意外送匿名评%
审。匿名评审把关学术部分，博士经过答辩，对这部分相对熟悉。反%
而是把关市场的出版社编辑，很难揣摩。这里面有很大的运气成分。%
最重要的是课题本身是否有潜在的读者市场。部分出版社愿意看我的%
书稿，大概就是稿子中的“耶稣会士”“天主教传播”和“东西文化”%
这些关键词自带话题性。后来，还有国内一知名出版社的编辑通过私%
人关系找到我，希望看我的书稿，可能也是这个原因。我遇到过一位%
研究中国上古音韵的欧洲学者，他就很无奈地告诉我，他领域的书稿%
很少有出版社愿意接手。%
\par
% 原书页码：379；段落
差不多2019年12月的时候，博睿编辑告诉我书稿已经外送，需%
要等待6—12个月的评审期。很快疫情开始，我在紧张的抗疫中渐渐%
忘了书稿的事。没想到2020年4月，博睿编辑直接发来了出版合同，%
合同的附件中，有两份审查报告。我打开几十页的报告，发现评审都%
是行家，优缺点都被他们找到了。二位评审均推荐出版。其中一位指%
出了不少问题，逐条列出，写了好几页。但在写完问题以后，他怕出%
版社编辑误解，专门说明这些问题只是为了让书稿更好，并不影响书%
稿整体的学术贡献，强烈建议立即出版。得到两位匿名评审的鼎力推%
荐，出版社直接给我发了合同，并且让我自己决定是否接受评审的修%
改建议。%
\par
% 原书页码：380；段落
我的英文书于 2021 年3 月出版上市，书名是 Jesuit Mission and%
\space%
Submission: Qing Rulership and the Fate of Christianity in China, 1644–1735.%
书稿出版，这个研究原本也就告一段落了。出乎意料的是，我陆续收%
到一些素不相识的学者来信，以及发在世界各大学术刊物上的8篇英%
文书评和2篇德文书评。一位给我写过书评的学者，在他的书评刊登%
后不久，通过一个学术网站发信联系我，询问我的中文姓名，还问我%
国内的师承。后来，一位学界前辈在他写的书评里，还专门附上了我%
的中文名。他们都在提醒我，我的书最直接的读者在中国。%
\par
% 原书页码：380–381；段落（跨页续段）
要把书变成中文版，是一个大工程。我博士答辩委员会中的周道%
本（JosephDauben）教授，也鼓励我出中文版。道本教授是研究欧洲%
数学史成名的，曾担任世界数学史协会主席。功成名就后，1980年代%
他偶然接触到中国数学，便从头开始学中文，一步一步学到可以阅读%
简牍的水平，成为西方世界研究中国古代数学的核心人物。当他跟我%
提议出中文版的时候，我找不到借口说写中文很难。2022年初，我试%
着把英文书的绪论翻译了出来，电邮发给了住在西昌师范学院的邓文%
树（1943—2023）先生。邓先生是经史子集、佛道音律全通的传统文%
人。只是为出身所累，一生隐居西南。邓先生看过绪论后，告诉我“句%
子很别扭”，要达到没有翻译味，需要大改。我自已也察觉到这一点，%
但无奈水平有限，很难去除翻译的痕迹。我遂放弃了翻译，变成了改%
写，因为我发现直接用中文书写其实还容易一些，这样写起来不受已%
有英文的约束，表达更自由。由于不是翻译，笔到之处，我也随着文%
思增减了许多内容。增减的标准中很重要的一条，就是要把故事写得%
“易读”和“好看”。这也是我之前学习“从论文到书”照顾市场的心%
得。学术味比较浓的绪论，我就将其分解，融入了各章节中。总体来%
说，中文版比英文版增加了约50\%的内容。大约也是提前想到了市场，%
第一位看稿的国内编辑就接受了我的书稿。%
\par
% 原书页码：381；段落
中文书稿最后出版上市，也算是我对国内的师友有了一个交代。%
首先要感谢的是当年为我写推荐信，帮我申请出国读博士的老师们。%
其中有主管我们班级的张翔教授，记得我当时骑自行车去张老师家中%
取推荐信，张老师端出水果，掰下一根香蕉递给我，对我说了很多鼓%
励的话。张老师早年留学日本，是日本史专家，对我出国留学的想法%
很支持。另一位是当年分管历史系本科生教育的章清老师。我向他提%
起要写推荐信，他很爽快地就答应了。%
\par
% 原书页码：381–382；段落（跨页续段）
另外两位给我写推荐信的老师已经作古。一位是王颋（1952—%
2018）教授。记得刚进大学，就是王颋老师给我们文史哲三系上大课，%
教中国古代史。他知识广博，每次来上课，总带着路上买的豆浆和%
蛋饼，站在讲台上匆匆吃完后，就根据当天的讲课内容，随手在黑%
板上画出讲课需要用到的历史地图，现在想来依然佩服不已。后来，%
我又选修了他的西北史地专业课。王颋老师读了我的期末论文《土%
尔扈特东归的“民族意义”别议》，不仅给了我A的成绩，还把这篇%
小文介绍到社科院办的学术网站上发表。为了出国留学，大四那一%
年，我除了准备英文考试以外，还希望能发表一些文章来帮助申请。%
当时我写了一篇《从〈世说新语〉看魏晋少年的早慧》，也是王颋教%
授帮我修改，后来投稿到中华书局的《书品》，还真被接收发表了。%
记得那时还收到了以汇款单形式邮寄给我的稿费。另一位给我写推%
荐信的是朱维铮（1936—2012）教授。朱维铮教授是历史系的前辈，%
并没有给本科生上课。我选修了他的学生廖梅老师开的“中国史学：%
史”。最后几堂课廖老师是请朱维铮教授来上的，他现场给我们写的%
文章打分。我当时得到了A-，这个成绩一直让我引以为傲，因为朱%
维铮老师从没有给出过A。后来经廖梅老师引荐，我找到朱维铮教授%
谈了出国读书的想法，正巧我要做的传教士选题和他对这段历史的%
理解有切合之处（这在本书第二章有呈现），朱维铮教授于是同意帮%
我出推荐信。还记得在签名前，他拿起笔，又放下，看着我说：“认%
真做点东西，不要只是出去混个文凭。”殷殷叮嘱，言犹在耳。%
\par
% 原书页码：382；段落
在决定写这部书后，我得到了多位大学同学的鼓励。特别是王晓%
刚、洪晓敏、夏炎和邹怡，他们通读了底稿，帮我梳理其中的错漏，%
在布局谋篇上为我提供建议。邹怡在读大学时就是系里学生刊物《史%
翼》的编辑，当年他就帮我改文章，使我在好几期《史翼》上都有发%
表。二十年后，再拿到他帮我改的书稿，有几分时光穿越之感。大学%
时的辅导员刘金华老师，以及当年和我一个寝室，现在复旦任教的徐%
建平同学，也帮我审读和修改了部分章节。2023年5月，在回上海参%
加毕业二十年同学聚会期间，刘老师和建平还帮我安排讲座，让我有%
机会报告自已的研究。感谢史翔和马晓艺同学，他们听说我有意出书，%
都对我鼓励有加，热心帮我联系出版事宜。%
\par
% 原书页码：382；段落
还要感谢人大清史所的胡祥雨。老胡是我在明尼苏达大学读书时%
的同学，当年我们一起合租房子，做了两年的室友。老胡是清史专家，%
是我遇到过历史知识最丰富的同辈学者。我在写作中，许多清史的问%
题都要请教他。我之前撰写英文专著，以及这次的中文著作，老胡都%
给了我很多指导。%
\par
% 原书页码：383；段落
出国快二十年了，一直忙于英文学习。这次是许多年来第一次正%
式用中文长篇写作。提笔之后才发现，自已英文进步不多，中文却退%
步很大。感谢本书责编胡晓凯女士帮我修改字句。胡编辑让我见识了%
年轻一代中文编辑的火眼金睛和字词功夫。%
\par
% 原书页码：383；段落
穿越历史，怀古伤今，有苦有乐。现实世界中身边都是些做金融%
的朋友，耳濡目染，学到了些市场翻腾的本事，也少了些柴米油盐的%
烦恼。初中都没毕业的父母经常在我耳边唠叨，祖上都是读书人，给%
我讲些“陋室空堂”“衰草枯杨”的前朝家事。“日月逝于上，体貌衰%
于下”，潜移默化中我也渐渐相信曹丕所谓“寄身于翰墨，见意于篇籍’%
是一种有意义的人生选择。好在妻子主外，我安于以 house husband自%
居，才有了这些年咬文嚼字、捻须苦吟的奢侈。%
\par
% 原书页码：383；段落
最后要感谢读者，能坚持读到这里。如果掩卷之后，能在互联网%
上惠赐您的感想，无论赞美还是批评，都是我的荣幸，我乐于拜读。%
\par
% 原书页码：383；段落
孙立天%
\par
% 原书页码：383；段落
2023年夏%
\par
% 原书页码：383；段落
于纽约郊外家中%
